\section{Interior phase inequalities}\label{int:section}

This section derives the local inequalities used in the cylinder argument.
Its input is the tangent construction of the preceding section, including
the information estimates \eqref{eq:vb} and \eqref{eq:fib}, the angular
profile \eqref{eq:ang}, and the projection calculus. We keep these inputs
explicit: an artificial quadratic branch does not acquire an angular
spreading hypothesis merely by having the same phase coordinates as an
equality branch.

The output is the list of admissible time--width directions in
Proposition~\ref{int:actions}. We first express changes of canonical entropy
through local phase-reading rates; the remaining arguments bound those
rates and the quadratic width defect.

\subsection{Tangent data and contact readings}

Fix a differentiability point $(a,v)$ in an interior patch for which the
replica construction is valid. Write
\[
 c=c',\qquad m=m(a,v),\qquad k_0=1-d,\qquad
 \delta=s^h,
\]
where $0<d\leq1$, $c>0$, and $0\leq m\leq1$. On a linear branch $c=b_1'$.
The calculations below have $g=0$. All auxiliary cutoffs and all finite
lists of depth queries are chosen before the inner scale limits. Errors
in information, divided by $\log(1/\delta)$, tend to zero before a
microlayer length tends to zero. Statements with an arbitrarily small
loss mean that this order is respected.

\begin{definition}[Admissible interior tangent data]\label{int:data}
The common root is $C_*$, the side is $Z$, and the phase $F=F(Z)$ is either
\[
 F=(X,B,P,Y,N,D)\quad\text{in the quadratic case},\qquad
 F=(X,B,Y,N,D)\quad\text{in the linear case}.
\]
Here $D$ denotes the normalized coordinate called $D'_0$ in the tangent
construction. Conditional on $Z$, finitely many observations are
independent with their prescribed common distribution. Their stored
marks are $\Xi_i=(\theta_i,r_i)$, together with $\mathfrak k_i$ in the
nonconstant quadratic case. Marks are stored to a sufficiently large
finite depth. Set $\mathcal D=(C_*,\boldsymbol\Xi)$.

We use the following properties of this construction.
\begin{enumerate}[label=(\roman*)]
\item The information estimates \eqref{eq:vb} and \eqref{eq:fib} hold
for every fixed finite set of phase, angular, and mark queries. In
particular the angular score at depth $p$ is $mp+o(1)$.
\item The nested phase and projection names satisfy the score calculus:
limiting scores have nonnegative Lipschitz increments; the chain rule,
conditioning inequalities, and comparisons by negligible lists hold
at actual samples; local mixed-prefix scores are linear at almost
every layer. The projection rule stated in the preceding section
applies with its specified finite-prefix Frostman hypotheses.
\item The time mark has dimension at least $d$. Given a joint mark
parent at depth $p>0$, the joint mark increment on $p\to p+t$ has
score at least $dt$, up to strict losses. For $c\geq1$ the time
increment itself has this lower bound. For a linear branch with
$c>1$, this remains true after conditioning on $r_{p+t}$, provided
$p+t<cp$. If $0<c<1$, the mixed mark profile supplies a number
$\nu\in[0,1]$ with
\[
 \begin{split}
 i_\delta(r_{p+t}\mid Z,\theta_{p+t},r_p)&=\nu t+o(t),\\
 i_\delta(\theta_{p+t}\mid Z,\theta_p,r_{p+t})
   &\geq (d-c\nu)_+t-o(t).
 \end{split}
\]
These assertions are made on the finite queries allowed by that
profile; $t$ is chosen small enough relative to $p$.
\item Resampling the side and phase conditionally on
$(C_*,\boldsymbol\Xi,Q_i)$ preserves the marked marginal. On a
resampling leg, with $A=\Delta X$, the increments obey
\eqref{eq:leg} to any fixed tested accuracy. The finite-walk
likelihood comparison supplied by \eqref{eq:fib} remains available
after conditioning on a common angular prefix when one is used.
\end{enumerate}
In the quadratic case, \emph{radial spreading} means that the
high alternative of \eqref{eq:rad} has $e_*\geq d$ almost surely for
the replicas. For conclusions involving $G_i$ when $c>1$, we also use
the stated uniform avoidance of affine lines. In the linear case,
\emph{crossing} means that both $|\theta_i-\theta_j|$ and
$|r_i-r_j|$ exceed $\delta^\chi$ with probability tending to one for
each fixed $\chi>0$.
\end{definition}

Subpower separation in this section always means the last type of
statement, with each positive exponent fixed before the inner limits.
It permits inversion with an arbitrarily small exponent loss. It does
not give a scale-independent lower bound.

Let $F_p$ be the isotropic phase name at depth $p$. On a layer
$p\to p+t$, $t\ll p$, use the contact forms
\begin{equation}\label{int:contact}
 H=dY-2X\,dB,\qquad J=dN-P\,dX,\qquad K=dD-P\,dB
\end{equation}
in the quadratic case; only $H$ is used in the linear case. These are
linear readings of increments, not differentials of the corresponding
products. For example, $J$ is not $d(N-PX)$.

For nested names, freeze $X,P$ at a fixed positive depth $p_0<p$.
After translation at the known parent $F_p$, changing this freeze
changes an increment by $O(\delta^{p+p_0})$. Genuine quadratic
remainders are $O(\delta^{2p})$. Both errors are negligible on a layer
with $t<p_0$ and $t<p$. Mark coefficients may be frozen in the same
way, at a fixed positive depth greater than $t$, keeping the known
translation at the parent.

Put $W_i=dB+\theta_i\,dX$. The readings needed below are
\begin{equation}\label{int:readings}
\begin{array}{lll}
 x_i=K+\theta_iJ-r_iH,&
 \pi_i=(W_i,H,x_i),&
 G_i=(L_i,J-\mathfrak k_iH,K-\beta_iH),\\[2pt]
 L_i=dP-2\beta_i\,dX+2\mathfrak k_i\,dB,&
 \beta_i=r_i-\theta_i\mathfrak k_i,&\text{quadratic},\\[4pt]
 x_i=dD+\theta_i\,dN-r_iW_i,&
 \pi_i=(W_i,H,x_i),&
 G_i=(dN-r_i\,dX,dD-r_i\,dB),\qquad\text{linear}.
\end{array}
\end{equation}
The quadratic $G_i$ is used on the equality branch with $c>1$.

Write $|U|(p)$ for the limiting random rate of the score of $U$ at
depth $p+t$, conditional on $F_p,\mathcal D$. Set
\[
 |U\mid V|=|(U,V)|-|V|,\qquad
 \alpha(p)=|F|(p),\qquad w(p)=\alpha(p)-m.
\]
The conditioning in $|U\mid V|$ includes the child reading of $V$.
Rates are defined simultaneously at typical layers for any finite or
countable collection of readings. Relative capacities and the
exponential bound for tiny cell probabilities identify their
expectations with entropy derivatives.

A bar denotes a common subsequential expected logarithmic Ces\`aro
mean at zero. Thus one averages against $dp/p$ over intervals whose
logarithmic lengths tend to infinity and whose lower endpoints tend
to zero. Constant dilations have the same mean. We take one
subsequence for the finite list of averages used in an argument.
Exchangeability gives equal expected means for different row indices.

\subsection{The cost of time and width}

\begin{proposition}[Differential cost]\label{int:differential-cost}
For admissible interior tangent data,
\begin{equation}
 \partial_a\mathcal H\geq d+
 \min(1,c)\bigl(\bar w+\overline{|x_i|}\bigr)
 +|1-c|
 \begin{cases}
  \overline{|\pi_i|},&c<1,\\
  \overline{|G_i|},&c>1,
 \end{cases}
 \tag{time-rate}\label{eq:time-rate}
\end{equation}
where the last term is omitted at $c=1$. In the linear case,
$\partial_v\mathcal H\leq\bar\alpha+1$.

In the quadratic case define, for $p\leq z\leq2p$, the Heisenberg
name $F_{p,z}$ by reading $X,B,P$ to depth $p$ and
\[
 Y-2[X]_pB,\qquad N-[P]_pX,\qquad D-[P]_pB
\]
to depth $z$. Let
\begin{equation}\label{int:width-defect}
 e(\lambda)=\frac{H_\delta(F_{2\lambda}\mid\mathcal D)
                -H_\delta(F_{\lambda,2\lambda}\mid\mathcal D)}{\lambda}
\end{equation}
in the limiting entropy law. Then
\begin{equation}
 \partial_v\mathcal H=2\bar\alpha-\bar e.
 \tag{width-rate}\label{eq:width-rate}
\end{equation}
\end{proposition}

\begin{proof}
Let $O$ be the actual child name at $(a+\lambda h,v)$. Given $C_*$,
its limiting entropy divided by $\lambda$ is $\partial_a\mathcal H$.
Even after giving a sufficiently fine phase name, the time bits cost
at least $d\lambda$: reveal $Z$ and apply the conditional time-mark
bound. The remaining entropy is bounded below by a sum of
\begin{equation}\label{int:layer-information}
 I_\delta(F_{p+t};O\mid C_*,F_p)
\end{equation}
over consecutive layers. The chain rule makes this a telescoping
sum of information about phase.

We will lower-bound the phase information by integrating the expected
values of the following local rates over the indicated intervals:
\[
\begin{array}{c|c}
 \text{phase-depth interval}&\text{local rate}\\ \hline
 (0,\min\{1,c\}\lambda)&w\\
 (c\lambda,\lambda),\quad c<1&|\pi_i|\\
 (\lambda,c\lambda),\quad c>1&|G_i|\\
 (\max\{1,c\}\lambda,(1+c)\lambda)&|x_i|.
\end{array}
\]
The middle interval is absent when $c=1$. The first rate comes
from recovering the phase after paying for its angle; the other
rates come from readings recoverable from the output.

For $p+t<\min(1,c)\lambda$, the output together with $C_*$ and
$X_{p+t}$ reads $F_{p+t}$ up to negligible lists. Its positions
give the two base coordinates at the child time; undoing its
displayed shear gives the common-frame coordinates and, in the
quadratic case, the residual derivative. The extra angular
increment has conditional cost at most $mt$ by
\eqref{eq:ang}. Consequently Equation~\eqref{int:layer-information}
pays the expected $w$-rate on this interval.

For the remaining intervals, we must identify the projected
readings and show that their coefficient information is paid for.
Let
\[
 \theta'=(t_{a+\lambda h}-t_a)/s^a,\qquad
 W'=B+\theta'X,\qquad V'=Y+\theta'X^2.
\]
The child-time error is $O(\delta^\lambda)$. The converted base
positions read $W',V'$ to depth $\lambda$. In common parent units,
the child shear difference has known constant offsets and
coefficients $\mathcal L,\mathcal R,\mathfrak b$ at time
$\theta'$. Its residual velocity and position are displayed to
depths $c\lambda$ and $(1+c)\lambda$. The derivative
\[
 \mathcal L+2\mathcal RX-2\mathfrak bW'
\]
approximates $P$ in the quadratic case, or $r_i$ in the linear
case, to depth $c\lambda$.

In a linear branch the curvature coefficients vanish, or are
negligible at every tested fixed power by the branch slack. For
constant quadratic cuts,
$\mathcal R=r_i+O(\delta^{c\lambda})$ and $\mathfrak b=0$.
For the quadratic equality branch with $c>1$,
\begin{equation}\label{int:output-coefficients}
 |\mathfrak b-\mathfrak k_i|+
 |\mathcal R-\theta'\mathfrak b-\beta_i|
 \lesssim\delta^{(c-1)\lambda}.
\end{equation}
This is the own-time coefficient comparison, extrapolated back
to normalized time zero after subtraction of the root trajectory.
All its nonconstant normalized coefficients are bounded; constant
offsets are retained exactly in the translations.

In the quadratic case the own-time algebra can be
checked without an implicit change of origin. Put
\[
 \mathfrak a=\mathcal R-\theta'\mathfrak b,\qquad
 \epsilon_0=P-\mathcal L-2\mathfrak aX+2\mathfrak bB.
\]
The derivative comparison gives
$|\epsilon_0|\lesssim\delta^{c\lambda}$.
With actual displayed coefficients held fixed, the
residual velocity has gradient
\[
 J-\mathfrak bH+\epsilon_0\,dX,
\]
and the residual position has gradient
\[
 K+\theta'J-\mathcal RH+\epsilon_0(dB+\theta'\,dX).
\]
Subtracting $\theta'$ times the former from the
latter gives $K-\mathfrak aH+\epsilon_0\,dB$.
These identities are evaluated at the actual phase
and follow directly by differentiating the shear
polynomials; the known constant offsets have zero
variation. At a parent representative the derivative
error is $O(\delta^p+\delta^{c\lambda})$.
Multiplication by the parent width gives
$O(\delta^{2p})+O(\delta^{p+c\lambda})$.
The second term is negligible on a microlayer
with $t<c\lambda$, and fits a binary layer
whenever $p\leq c\lambda$.
The derivative display itself
has gradient $dP-2\mathfrak a\,dX+2\mathfrak b\,dB$.
Equation~\eqref{int:output-coefficients} therefore
gives the three claimed $G_i$ readings with the
stated precision.

In the linear case the shear is linear, with
$|r_i-\mathcal L|\lesssim\delta^{c\lambda}$.
The displayed residual velocity and the backflowed
residual position have respective gradients
\[
 dN-\mathcal L\,dX,
 \qquad dD-\mathcal L\,dB.
\]
They differ from the two components of $G_i$
by $(r_i-\mathcal L)dX$ and
$(r_i-\mathcal L)dB$. Before backflow the
position gradient is
$dD+\theta'\,dN-\mathcal L(dB+\theta'\,dX)$,
which gives $x_i$ with the additional
$O(\delta^\lambda)$ time-coefficient error.

Center each displayed quantity by its value at $F_p$, evaluated
with its actual displayed coefficients and offsets. Taylor's
formula has error $O(\delta^{2p})$. The residual-position
gradient converges, with a positive power error on an interior
layer, to $x_i$. Below depth $\lambda$, the base-position
gradients additionally give $W_i,H$. If $c>1$, subtract
$\theta'$ times the displayed velocity from the position.
In the quadratic case the two resulting gradients give
$K-\beta_iH$ and $J-\mathfrak k_iH$, and the derivative equation
gives $L_i$. In the linear case they give the two components
of $G_i$. Thus the intermediate interval gives $\pi_i$ when
$c<1$ and $G_i$ when $c>1$, while the final interval always
gives $x_i$.

For completeness, the projected entropy is indeed paid by
Equation~\eqref{int:layer-information}, even though its coefficients
may depend on phase. Freeze the needed angular and curvature
coefficients into a name $\Omega$ at a depth $q$ with
\[
 t<q<
 \begin{cases}
 \min(p,\lambda,c\lambda),
    &\text{linear or constant quadratic cuts},\\
 \min(p,\lambda,(c-1)\lambda),
    &\text{nonconstant quadratic cuts, }c>1.
 \end{cases}
\]
Shrink the microlayer first so that such a fixed $q$
exists. On a compact interior interval of depths,
the preceding output comparisons recover $\Omega$ from
$O,C_*,F_p$ up to negligible lists. The centered reading $U$ is
also determined, up to such lists, by both
$(F_{p+t},C_*,F_p,\Omega)$ and $(O,C_*,F_p,\Omega)$.
For finite tested names $A,B,\Omega,U$ and arbitrary
conditioning data $C$, the chain
rule and the two reconstruction bounds give
\[
 \begin{split}
 I_\delta(A;B\mid C)\geq {}&H_\delta(U\mid C,\Omega)
          -H_\delta(\Omega\mid B,C)\\
 &-H_\delta(U\mid A,C,\Omega)
          -H_\delta(U\mid B,C,\Omega).
 \end{split}
\]
To obtain this inequality, first adjoin $\Omega$ to $B$,
at cost at most $H_\delta(\Omega\mid B,C)$, retain the
conditional information given $\Omega$, and apply the
chain rule to the two reconstructions of $U$.
Use $A=F_{p+t}$, $B=O$, and $C=(C_*,F_p)$.
All three error terms tend to zero at the fixed layer,
before division by $t$. Conditioning further on the stored marks lowers this
entropy. Translation at the parent and the coefficient-freezing
comparison then identify the relevant contact-reading rate.

Tile compact subintervals of the three depth ranges by smaller
layers. Bounded convergence for the score derivatives and
nonnegativity for the information sums permit passage to their
integrals. Truncate near the endpoints, then remove the
truncations using the uniform relative capacities. Finally
average in $\lambda$ logarithmically. Dilation invariance of
this mean gives Equation~\eqref{eq:time-rate}.

For width variation, increasing $v$ by $\lambda h$ in a linear
branch reads at most $F_\lambda$ and one additional
$\lambda$-depth scalar in vertical position. This proves the
linear upper bound. In a quadratic branch the same change
refines $X,B,P$ by $\lambda$ and the three recentered
coordinates by $2\lambda$. In changing the linear shear,
the new binned coefficient differs from the true normalized
$P$ by $O(\delta^\lambda)$. After translation, its effect on
unresolved horizontal widths is $O(\delta^{2\lambda})$.
Hence the new width name and $F_{\lambda,2\lambda}$ determine
one another up to bounded lists, given $C_*$.
Equation~\eqref{eq:vb} allows marked conditioning in these
finite-depth entropy comparisons. Since $F_0$ has zero
normalized entropy, logarithmic averaging of
Equation~\eqref{int:width-defect} proves
Equation~\eqref{eq:width-rate}.
\end{proof}

\subsection{Angular modes and projection amplification}

\begin{lemma}[Angular and commutator modes]\label{int:modes-geometric}
The angular fiber direction is
\begin{equation}\label{int:angular-mode}
 \ell_i=
 \begin{cases}
 (1,-\theta_i,2r_i;0,0,0),
       &(dX,dB,dP;H,J,K)\quad\text{quadratic},\\
 (1,-\theta_i,r_i,-\theta_ir_i;0),
       &(dX,dB,dN,dD;H)\quad\text{linear}.
 \end{cases}
\end{equation}
Any nondegenerate scalar component in this direction has
conditional rate at least $m$ after adjoining finitely many
simultaneous readings that annihilate $\ell_i$. A conditional
angular-coordinate rate is at most $m$.

In the quadratic case, if $m>0$, the center also has the mode
\begin{equation}\label{int:commutator-mode}
 b_{ij}=(\theta_i-\theta_j,\ r_i-r_j,\
                    \theta_ir_j-\theta_jr_i)
       \quad\text{in }(H,J,K).
\end{equation}
A nondegenerate component of this mode has residual rate at
least $m$ after conditioning on center readings that annihilate it.
Both $x_i$ and $x_j$ annihilate $b_{ij}$.
\end{lemma}

\begin{proof}
Adjoin $Q_i,X_p$ to the conditioning. Along this fiber the
phase is determined by $X$ to every fixed tested accuracy.
Differentiating the fiber equations gives
Equation~\eqref{int:angular-mode}. Their quadratic remainder
on an angular parent cell is $O(\delta^{2p})$. Thus $F_p$
and the annihilating child readings have negligible list
cost with the enlarged conditioning. A nondegenerate
component recovers $X_{p+t}$ up to negligible lists. For a
subpower lower bound on its derivative, fix the separation
exponent much smaller than the strict loss in the test,
apply one-dimensional inversion, and then let that exponent
decrease. Equation~\eqref{eq:fib} gives the residual rate
$m$. The reverse upper bound for angular rates follows
from \eqref{eq:ang} and conditioning.

For the center mode, walk successively along fibers $i,j$,
with endpoints $F_0,F_1,F_2$ and angular coordinates
$X_0,X_1,X_2$. Adjoin the starting side, very fine $X_2$
bits, and $X_{1,p}$. The finite-walk likelihood comparison
and \eqref{eq:fib} leave rate $m$ for $X_{1,p+t}$. The
separation $|X_1-X_2|\geq\delta^\chi$ holds with probability
tending to one for every fixed $\chi>0$: conditional on the
previous history, a new angular prefix has positive score
at each positive depth, so it cannot lie in a bounded list
of prescribed bins with nonvanishing probability.

Apply \eqref{eq:leg} on the two legs and differentiate with
$X_2$ fixed. The resulting center derivative at $F_2$ is
\[
 2(X_2-X_1)b_{ij}\,dX_1.
\]
In particular its last component is
$2(X_2-X_1)(\theta_ir_j-\theta_jr_i)dX_1$.
The same parent-cell Taylor and inversion argument applies.
The endpoint has the original marked marginal, so the
result is a statement about its original conditional rates.
Substitution in the formula for $x_i$ or $x_j$ verifies
the last assertion.
\end{proof}

\begin{lemma}[Projection amplification]\label{int:boost}
Let a planar pair have joint local rate $S$, possibly after
conditioning on common finer readings. Suppose that a scalar
projection of rate $x$ leaves a residual rate at least $M>0$.
Suppose also that its direction has Frostman exponent $d$ on
the finite prefixes required by the projection rule, and that
its full varying record has negligible information with the
planar source. Then
\[
 x\geq\min(1,d+M).
\]
\end{lemma}

\begin{proof}
The chain rule gives $S\geq x+M$. The projection rule gives
$x\geq\min(1,S,(S+d)/2)$. If $x<1$, then $S>x$, so the
only remaining possibility is $2x\geq S+d\geq x+M+d$.
This gives $x\geq M+d$.

For random rates these deductions are applied to a
positive-probability set with a strict deficit. Bin the
rates to sufficiently short fixed intervals, use mixed-prefix
linearity to obtain the finite-prefix source bounds, and
use popular projection cells to obtain the output lists.
The information error vanishes at this fixed layer. The
projection rule excludes the deficit. This procedure also
allows the high-set radial witnesses and the low-set
pin-product witnesses of the preceding section; it never
conditions on an unspecified subpower-probability event.
\end{proof}

\begin{proposition}[Basic phase rates]\label{int:basic-rates}
Put $l_* = \min(1,d+m)$. For quadratic data with $m>0$ and
radial spreading, write $C_z=(H,J,K)$. Then
\begin{equation}
 |W_i\mid C_z|\geq l_*,\qquad
 |C_z|-|x_i|\geq m+l_*,\qquad
 \alpha-|C_z|\geq2m+l_*.
 \tag{modes}\label{eq:modes}
\end{equation}
When the quadratic $G_i$ is used and line avoidance holds,
\begin{equation}
 |G_i|\geq2m+|x_i|,\qquad
 2\E|G_i|\geq\E\alpha-1.
 \tag{Gq}\label{eq:Gq}
\end{equation}
For linear data with $m>0$ and crossing,
\begin{equation}
 |W_i\mid H|\geq l_*,\qquad |x_i|\geq l_*,\qquad
 |G_i|\geq m+l_*,\qquad
 2\E|G_i|\geq\E\alpha-1.
 \tag{modes-lin}\label{eq:modes-lin}
\end{equation}
For linear data, without a positivity assumption on $m$,
\begin{equation}
 w-|\pi_i|\leq1.
 \tag{miss}\label{eq:miss}
\end{equation}
\end{proposition}

\begin{proof}
In all uses of Lemma~\ref{int:boost}, the planar source and
any common child conditioning depend only on phase and the
held records. Equation~\eqref{eq:vb}, followed by the
conditional chain rule, permits removal of the varying
record. Its directional sublaw bounds, initially conditional
on $Z$ and the held records, integrate to the tested
conditioning. This verifies the information requirement
without asserting independence after a witness restriction.

For $W_i$, use the source $(X,B)$, conditional on the center
in the quadratic case and on $H$ in the linear case.
The direction parameter is $\theta_i$, with exponent $d$,
and $\ell_i$ gives residual $m$.

For the quadratic center inequality, condition on $x_i$
to child depth and use $(H,J)$. It determines all center
coordinates with $x_i$. Vary $j$ and project with kernel
the image of $b_{ij}$. The direction is a nonsingular
linear transform of the joining-line direction in mark
space. Radial spreading and Lemma~\ref{int:modes-geometric}
give a projected rate at least $l_*$ and residual $m$.
Their sum gives $|C_z|-|x_i|\geq m+l_*$.

For the remaining horizontal rate, condition on the whole
center and quotient $(dX,dB,dP)$ by $\ell_i$. One may use
the quotient coordinates $(W_i,dP-2r_i\,dX)$. The missing
coordinate has rate at least $m$. On the quotient, vary
$j$ and project with kernel the image of $\ell_j$. Its
direction is another nonsingular transform of the joining
line, and its residual is $m$. Thus the quotient rate is
at least $m+l_*$, proving the third bound in
\eqref{eq:modes}. Time separation suffices for the
nondegenerate coordinate comparisons.

Next, $x_i$ is a combination of the two center components
of quadratic $G_i$. Given $x_i$, the form $J-\mathfrak k_iH$
sees $b_{ij}$. Given the center, $L_i$ sees $\ell_j$.
Both tests are nondegenerate because
\[
 |r_j-r_i-\mathfrak k_i(\theta_j-\theta_i)|
\]
has subpower separation by line avoidance. They contribute
two residual copies of $m$. The pair $G_i,G_j$ determines
all three center coordinates and two independent horizontal
combinations. Indeed the two coefficient pairs
$(\beta_i,\mathfrak k_i)$ and $(\beta_j,\mathfrak k_j)$ have
subpower separation by the same line test. One scalar
coordinate remains, with rate at most one. Subadditivity
and exchangeability prove the second bound in
\eqref{eq:Gq}.

For the linear $x_i$ bound, use the planar source
$(dX,x_i)$ and put
$f=(\theta_i-\theta_j)(r_j-r_i)$. Both projections
$x_i$ and $x_i-f\,dX$ have residual rate exactly $m$:
$\ell_i$ and $\ell_j$, respectively, give the lower
bound, and the angular upper bound gives equality.
Consequently their projected rates agree. Vary $j$ in
the second projection and use the pin-product direction
alternatives with Lemma~\ref{int:boost}. This proves
$|x_i|\geq l_*$. For $G_i$, its image of $\ell_j$ is
nondegenerate by crossing and has direction $\theta_j$.
The same argument gives a projected rate $l_*$ and
residual $m$. Two $G$ readings determine every coordinate
except $Y$, so the last expected inequality follows.

Finally $(\pi_i,X)$ misses only one additional scalar
coordinate by rank. In this comparison all inversions
have bounded coefficients. Its missing angular rate is
at most $m$, and the other scalar costs at most one.
Thus $\alpha\leq|\pi_i|+m+1$, which is
\eqref{eq:miss}.
\end{proof}

\begin{lemma}[A many-leg vertical rate]\label{int:many-leg}
For linear admissible tangent data with $m>0$,
\begin{equation}
 |H|\geq l_*.
 \tag{H-walk}\label{eq:H-walk}
\end{equation}
For each strict-error application a sufficiently large
finite number of replicas suffices.
\end{lemma}

\begin{proof}
Walk $2n$ distinct legs. Adjoin the starting side, very
fine even angular coordinates, including $X_{2n}$, and
the odd angular parents at depth $p$. These data predict
the endpoint parent to negligible lists. Summing the
$B,Y$ formulas in \eqref{eq:leg}, and translating at that
parent, makes the endpoint $H$ reading equivalent to
\[
 S_n=\sum_{k=1}^n(\theta_{2k}-\theta_{2k-1})
       (X_{2k-1}^2-2X_{2n}X_{2k-1}).
\]
The coefficient of $H$ may first be frozen at the known
fine endpoint or at the prescribed fixed positive prefix;
the increment comparison gives the same rate.

Let $s_k$ be the rate of the partial sum with these data,
and set $s_0=0$. The source consisting of $S_{k-1}$ and
the new odd quadratic scalar has joint prefix rate at
least $s_{k-1}+m$. To see the second contribution, reveal
the other fine endpoint bits, which predict the previous
sum. The derivative of the new scalar is
$2(X_{2k-1}-X_{2n})$, with subpower separation, and
\eqref{eq:fib} gives its remaining angular rate $m$.

For the projection, omit the current pair of records.
The difference of their time marks has exponent $d$.
Here is the required conditional comparison. If $X_j^*$
are fixed finite-depth angular names fine enough for
all these tests, the walk rule gives
\begin{equation}\label{int:walk-kl}
 D_{\mathrm{KL}}\left(
  \Prob_{Z_0,\boldsymbol\Xi,(X_j^*)_j}
  \,\middle\|\,
  \Prob_{Z_0,\boldsymbol\Xi}
       \otimes_{C_*}\prod_j\Prob_{X^*\mid C_*}
             \right)=o(\log(1/\delta)).
\end{equation}
This follows by the conditional chain rule in chronological
order: the next endpoint is sampled from its conditional
fiber kernel, the marked marginal is stationary, and its
angular information cost is \eqref{eq:fib}. The initial
angular coordinate is included in $Z_0$ and need not be
independent of it.

The reference law in Equation~\eqref{int:walk-kl}
factorizes the omitted record and the source angular
bits conditionally on the remaining data. The source
parents are predicted to negligible lists by the odd
parents and the retained data. Conditional chain rules
give the same negligible information budget after
these parents are revealed. Conditional likelihoods
at actual samples agree up to vanishing exponents.
Trim excessive likelihood ratios and time-ball failures
with any fixed small power slack. The unchanged
start-side marked marginal supplies the unnormalized
direction bounds, which integrate to the source
conditioning. This verifies the projection-rule
hypotheses at the actual law; no assertion merely
typical for the reference law is transferred without
a likelihood bound.

The projection rule therefore gives
\[
 s_k\geq\min\left(1,s_{k-1}+m,
                         \frac{s_{k-1}+m+d}{2}\right).
\]
Iterating this monotone map from zero converges to
$\min(1,m+d)$: below that value each of its three
entries exceeds the current value, and its least
positive fixed point is $\min(1,m+d)$. Adding the
auxiliary data can only lower the original endpoint
score, whose parent was already predicted. Given any
strict desired error, choose a finite $n$ making the
iterate sufficiently close to the limit, and only
then choose the common replica experiment and cutoffs.
\end{proof}

\subsection{Orientation motions}

The preceding estimates varied the angular coordinate with a mark held
fixed. We now vary the mark in the equations relating two phases with
shared fiber data. Holding one phase and the target angle fixed determines
the resulting motion of the target phase. These motions make the time and
orientation innovations available as further phase readings.

\begin{lemma}[Mark-motion comparison]\label{int:mark-motion}
Assume $m>0$. Fix a record $j$ and resample one or two
sides $R_s$ independently from the conditional kernel
with the same $Q_j,\mathcal D$ as a target $F$. Put
$A_s=X_F-X_{R_s}$. Each $A_s$, and $A_1-A_2$ when two
draws are used, has subpower separation. With $R_s,X_F$
held to sufficiently fine tested precision, variation
of the mark $\xi_j=(\theta_j,r_j)$ has differential
\begin{equation}
\begin{array}{ll}
 dX=0,\quad dB=-A\,d\theta_j,\quad H=A^2\,d\theta_j,
      &\text{both cases},\\
 dP=2A\,dr_j,\quad J=A^2\,dr_j,\quad
 K=A^2(r_j\,d\theta_j-\theta_j\,dr_j),
      &\text{quadratic},\\
 dN=A\,dr_j,\quad
 dD=-r_jA\,d\theta_j-\theta_jA\,dr_j,
      &\text{linear}.
\end{array}
\tag{trans}\label{eq:trans}
\end{equation}
Here $A=A_s$, and the forms are frozen at the target.
In particular $x_j$ annihilates the full two-component
motion. These motions give conditional lower-score
tests with the mark innovation bounds of
Definition~\ref{int:data}, when the held readings
annihilate the tested innovation.
\end{lemma}

\begin{proof}
The target and the resampled sides are exchangeable
conditionally on the shared fiber data. Their angular
information bounds give the stated separations.
Holding $R_s$ and $X_F$ fixes $A$. In the quadratic
case the leg formulas give
\[
 \begin{aligned}
 B_F&=B_R-\theta_jA,&
 Y_F&=Y_R-\theta_j(2X_RA+A^2),\\
 P_F&=P_R+2r_jA,&
 N_F&=N_R+P_RA+r_jA^2,\\
 D_F&=D_R-\theta_j(P_RA+r_jA^2).
 \end{aligned}
\]
Differentiating and substituting $X_F=X_R+A$ and
$P_F=P_R+2r_jA$ into the contact forms gives
Equation~\eqref{eq:trans}. In the linear case use
$N_F=N_R+r_jA$ and $D_F=D_R-\theta_jr_jA$.
Direct substitution proves the assertion about $x_j$.

The conditioning order is essential. First freeze every
coefficient involving the moving mark at a depth below
$p$ and above $t$. With the full record known, translate
at $F_p$ to compare the original and frozen readings.
Equation~\eqref{eq:vb} then removes the remaining fine
record from the list test. Only after this removal do
we adjoin $Z_{R_s}$ and the motion data. At a mark parent
of depth $p$, these data predict $F_p$ and every held
annihilating child reading up to negligible lists.
An informative component recovers the scalar innovation
by inversion, with arbitrarily small power loss. Its
Taylor error is $O(\delta^{2p})$. Holding the other
mark scalar to depth $p+t$, when permitted by its
innovation law, adds only an error of that order
of accuracy.

Auxiliary angular bits do not invalidate removal of
the moving record. For example, for two draws and
fixed sufficiently fine names,
\begin{equation}\label{int:motion-independence}
 I_\delta((X_{R_1}^*,X_{R_2}^*);Z_F,\boldsymbol\Xi
                  \mid C_*)=o(1).
\end{equation}
Indeed let $\mathcal B=(Q_j,\mathcal D)$ and
$Y_s=X_{R_s}^*$. The variables $Y_1,Y_2$ are
conditionally independent given $\mathcal B$,
and their joint conditional law depends on
$(Z_F,\boldsymbol\Xi)$ only through
$\mathcal B$. Since $\mathcal B$ is determined
by these latter data, the conditional chain rule
gives
\[
 \begin{split}
 I_\delta((Y_1,Y_2);Z_F,\boldsymbol\Xi\mid C_*)
 &=I_\delta((Y_1,Y_2);\mathcal B\mid C_*)\\
 &\leq\sum_{s=1}^2 I_\delta(Y_s;\mathcal B\mid C_*)
 =o(1).
 \end{split}
\]
The inequality subtracts the nonnegative term
$I_\delta(Y_1;Y_2\mid C_*)$; the last equality
is \eqref{eq:fib} for the stationary marginals.
The target
bits $X_F^*$ may be included in the finite phase
budget in \eqref{eq:vb}. At $R_s$, its own angular
bits are known, and the other angular bits satisfy
the corresponding version of
Equation~\eqref{int:motion-independence}. Finally,
$(Z_{R_s},\boldsymbol\Xi)$ has the original marked
marginal, so the replicas are conditionally independent
there and the stated mark innovation law applies.
Conditional chain rules permit the specified mark
parents and child scalar bins. This proves the score
comparison without asserting independence conditional
on all the exact motion data simultaneously.
\end{proof}

\begin{proposition}[Orientation rates]\label{int:orientation-rates}
For quadratic data with $m>0$ and radial spreading,
\begin{equation}
 |x_i|\geq d.
 \tag{cross-q}\label{eq:cross-q}
\end{equation}
If also $c<1$, set $a_*=(d-c\nu)_+$ and
$n_* =\min(1,d+\nu)$. Then
\begin{equation}
 |x_i|\geq n_*,\qquad |H\mid x_i|\geq a_*,\qquad
 |C_z|-|x_i|\geq a_*+\nu,\qquad
 |\pi_i|\geq l_*+a_*+|x_i|.
 \tag{mix-q}\label{eq:mix-q}
\end{equation}
For linear data with crossing, $m>0$, $c<1$, and $\nu>0$,
\begin{equation}
 |x_i|\geq n_*,\qquad |\pi_i|\geq2l_*+n_*,\qquad
 w-|\pi_i|\geq\nu.
 \tag{mix-l}\label{eq:mix-l}
\end{equation}
The many-leg term in this assertion is obtained with
an arbitrary strict small loss using finitely many replicas.
\end{proposition}

\begin{proof}
For the quadratic $x_i$ bound, move a record $j\ne i$.
Its gradient in mark space is
\[
 A^2(r_j-r_i,\theta_i-\theta_j).
\]
Keep $Z_R$, the moving mark parent $\xi_{j,p}$, the
needed target angular bits, and other records except
$i,j$. Use $i$ as the varying projection record.
The input mark increment has dimension at least $d$.
By the parent-representative version of radial
spreading, the gradient directions over $i$ have
unnormalized Frostman exponent $d$ on the high
witness relation. Swapping and signing coordinates
does not change that assertion. Pair separation
makes the gradient magnitude nondegenerate up to
subpower loss.

To verify the output and independence hypotheses,
start from a proposed low-score list for $x_i$
conditional on $F_p,\mathcal D_{-j}$. The removal of
record $j$ is allowed by \eqref{eq:vb}. Given its
parent and the motion data, $F_p$ has negligible
lists. Taylor expansion at the phase predicted by
the mark parent, followed by dilation by
$\delta^{-p}$, turns this into a list for the stated
linear projection. The law of the angular bits has
negligible likelihood cost relative to retaining
$(Z_R,\boldsymbol\Xi)$ and sampling those bits
independently given $C_*$. Under this comparison
kernel the planar source and the projection record
are independent given the test data. Trim likelihood
and direction failures with fixed small power slack;
the marked marginal is unchanged. Thus the source
and direction dimensions are both at least $d$,
and the projection rule gives $|x_i|\geq d$.
Precisely, Lemma~\ref{names:parent-direction-transfer}
applies with side $Z_R$, auxiliary name $X_F^*$,
$S=C_*$, moving mark $j$ at parent depth $b=p$,
and pin record $i$. It supplies both the
unnormalized parent-direction bound and its
conditional likelihood transfer, without dividing
the pin bound by the moving parent's probability.

For the mixed quadratic bounds, move $i$ itself.
Hold $r_i$ to child depth. The $H$ component reads
the time innovation, of rate at least $a_*$, and
$x_i$ annihilates it. Next hold $\theta_i$ to child
depth. The $J$ component reads the $r_i$ innovation,
of rate $\nu$, while both $H$ and $x_i$ may be held.
The chain rule gives the two center bounds.
If $\nu>0$, use the planar pair $(J,K)$ conditional
on $H$ and the projection $K+\theta_iJ$. Holding
$\theta_i$ to child depth, this projection and $H$
annihilate the $r_i$ innovation, while $J$ reads it.
Remove the full $i$ record from the planar source
by \eqref{eq:vb}, then use the time direction and
Lemma~\ref{int:boost} with residual $\nu$.
The resulting projected rate is $n_*$. Restoring
the record and $H$ identifies it with $x_i$.
When $\nu=0$ use \eqref{eq:cross-q}. Finally
$W_i$, even conditional on the whole center,
costs at least $l_*$ by \eqref{eq:modes}.

In the linear case condition on $(X,B,H)$ to child
depth. The planar source is $(dN,dD)$ and the
projection is $dD+\theta_i\,dN$. With $\theta_i$
held to child depth, the $r_i$ innovation has rate
$\nu$, is annihilated by this projection, and is
read by $dN$. Lemma~\ref{int:boost} gives rate
$n_*$, also for $x_i$ with this finer conditioning.
Add $|W_i\mid H|\geq l_*$ and
$|H|\geq l_*$ from \eqref{eq:H-walk}. This proves
the first two claims. Given $\pi_i$, the angular
mode gives rate $m$ for $X$; after those readings
are held, $dN$ still reads the $r_i$ innovation
with residual $\nu$. Consequently
$\alpha\geq|\pi_i|+m+\nu$.
\end{proof}

\begin{proposition}[Two time innovations]\label{int:double}
For linear admissible data with $m>0$ and crossing,
\begin{equation}
 \begin{array}{ll}
 |\pi_i|\geq2d+l_* ,&c<1,\ \nu=0,\\
 w-|x_i|\geq2d,&c=1,\\
 w-|G_i|\geq2d,&c>1.
 \end{array}
 \tag{double}\label{eq:double}
\end{equation}
\end{proposition}

\begin{proof}
In the first case move $j\ne i$, holding $r_j$ to
child depth. The time innovation still has rate at
least $d$. The reading
\[
 V=dD+\theta_i\,dN-r_jW_i
\]
annihilates that motion. It is a projection of
$G_j$ with direction $\theta_i$. The image of
$\ell_i$ in $G_j$ is nondegenerate by crossing
and is annihilated by $V$. Lemma~\ref{int:boost}
therefore gives $|V|\geq l_*$.

For $c=1$ move $i$ and set $V=x_i$. Do not hold
the refined $r_i$; the joint-parent time innovation
still has rate $d$, and $x_i$ annihilates the full
motion. For $c>1$, move $i$, set $V=G_i$, and hold
$r_i$ to child depth. The allowed depth condition
$p+t<cp$ gives the time rate $d$, and $G_i$
annihilates this remaining motion.

In all cases use two conditional draws $R_1,R_2$
as in Lemma~\ref{int:mark-motion}. The motion of
$(W_i,H)$ has components $(-A_s,A_s^2)d\theta$.
After adding the auxiliary angular bits, first
test $H+A_1W_i$ with the second draw. Its
coefficient is $A_2(A_2-A_1)$, with subpower
separation, and $V$ is held. Next hold both $V$
and $H+A_1W_i$, and test $W_i$ with the first
draw. The held combination annihilates this
motion and the coefficient of $W_i$ is $-A_1$.
The chain rule gives
$|(W_i,H)\mid V|\geq2d$. The conditioning
comparison in Lemma~\ref{int:mark-motion}
justifies each test before dropping the
auxiliary bits.

In the first case $(W_i,H,V)$ and $\pi_i$ are
equivalent with the marked data, since their
last components differ by $(r_i-r_j)W_i$.
In the other two cases the held readings also
annihilate $\ell_i$, leaving angular rate $m$.
These observations prove all three inequalities.
\end{proof}

\subsection{Quadratic width and binary time layers}

Two estimates remain before we can collect the admissible directions.
The first bounds the entropy cost of refining the horizontal width in a
quadratic phase. The second uses binary time layers to improve the time
bound at large quadratic cost.

\begin{proposition}[Horizontal width]\label{int:horizontal-width}
For quadratic admissible tangent data,
$\bar e\geq2m$. If $m>0$ and radial spreading holds,
\begin{equation}
 \bar e\geq2m+d.
 \tag{e-q}\label{eq:e-q}
\end{equation}
\end{proposition}

\begin{proof}
Fix $p<z<2p$ and refine the horizontal depth from
$p$ to $p+t$, leaving $z$ fixed. Given $F_{p,z}$,
this increment is equivalent to the new horizontal
bits: recentering contributes only
$O(\delta^{2p})$. For every angular mode,
$Q_i,X_p,\mathcal D$ predict $F_{p,z}$ up to
negligible lists. Indeed the contact-center
derivative on the fiber is $O(\delta^p)$ on
the parent, and its variation is
$O(\delta^{2p})$.

Two time-separated angular modes are linearly
independent. First test a horizontal scalar
annihilating $\ell_i$ and seeing $\ell_j$;
then test a scalar seeing $\ell_i$. These
tests remain valid after holding the scalar
annihilating both modes. Lemma~\ref{int:modes-geometric}
therefore gives horizontal increment rate
at least $2m$. Integrating in $p$ from
$\lambda$ to $2\lambda$ with $z=2\lambda$
fixed yields $e(\lambda)\geq2m$ in the
averaged sense. Boundary endpoints are
handled by the Lipschitz continuity of
the array on the wedge.

For the additional rate, a horizontal scalar
annihilating both modes is
\begin{equation}\label{int:normal-horizontal}
 U=(\theta_i-\theta_j)(dP-2r_i\,dX)
       -2(r_j-r_i)(dB+\theta_i\,dX).
\end{equation}
Set $q_0=z-p$ and $u=p-q_0>0$. Resample $R$
on the $i$ fiber while sharing $X_{q_0}$.
Then $A=X_F-X_R$ is $O(\delta^{q_0})$ and,
for every fixed $\chi>0$, at least
$\delta^{q_0+\chi}$ in magnitude with
probability tending to one. This follows
from the conditional version of
\eqref{eq:fib}. The marked marginal is
stationary, and the target angular bits
beyond the common prefix retain negligible
information with $Z_R,\boldsymbol\Xi$
given $C_*,X_{q_0}$. To see the exact comparison,
set $\mathcal S=(C_*,X_{q_0})$ and
$\mathcal B=(Q_i,\mathcal D,X_{q_0})$.
The resampling kernel is the conditional side
law given $\mathcal B$. The conditional Markov
property and \eqref{eq:fib} give
\[
 \begin{split}
 I_\delta(X_F^*;Z_R,\boldsymbol\Xi\mid\mathcal S)
 &\leq I_\delta(X_F^*;Q_i,\mathcal D
                             \mid C_*,X_{q_0})\\
 &\leq I_\delta(X_F^*;Q_i,\mathcal D\mid C_*)=o(1).
 \end{split}
\]
The second inequality uses that $X_{q_0}$ is a
prefix of $X_F^*$. Thus the comparison law keeps
$(Z_R,\boldsymbol\Xi)$ unchanged and samples
$X_F^*$ independently with its conditional law
given $\mathcal S$. The common prefix is already
determined by $Z_R$.

Start the moving mark at depth $u$. Given
$Z_R$, fine $X_F$ bits, and $\xi_{i,u}$,
the name $F_{p,z}$ has negligible lists.
Horizontal uncertainty is
$O(A\delta^u)=O(\delta^p)$. Guess the
boundedly many corresponding horizontal
bins. With these bins fixed, the recentered
uncertainty is
\[
 O(A^2\delta^u)+O(\delta^{2p})
        =O(\delta^z).
\]
The mixed quadratic mark remainder is
$O(A^2\delta^{2u})$, which is smaller.

Freeze the $i$ coefficients in
Equation~\eqref{int:normal-horizontal}
at their $u$-parent representatives.
Indeed its expanded form is
\[
 U=(\theta_i-\theta_j)dP-2(r_j-r_i)dB
                  -2(\theta_ir_j-\theta_jr_i)dX.
\]
Each coefficient changes by $O(\delta^u)$.
On a horizontal parent its centered increment
has size $O(\delta^p)$, so the change is
$O(\delta^{p+u})$, negligible at depth $p+t$
when $t<u$. Absolute shifts at the parent
are retained during this comparison with the
full record. Remove the
remaining fine $i$ mark by
\eqref{eq:vb} before adjoining motion
data. With $X_F$ fixed, the resulting
mark-plane projection is
\[
 2A\bigl((r_j-r_{i,u})\,d\theta_i
          +(\theta_{i,u}-\theta_j)\,dr_i\bigr).
\]
After dilation of the mark parent,
a low-rate list at depth $p+t$ becomes
a projection list at depth $t$, up to
arbitrarily small power losses. Its
source has joint dimension at least $d$.
More explicitly, the common test data are
$Z_R,X_F^*,\xi_{i,u}$ and the held records
other than $i,j$. The source is the child
increment of $\xi_i$ and the projection
record is $j$. Under the comparison law
they are independent given these data:
the marks are independent given $Z_R$,
and $X_F^*$ is independent of them given
$\mathcal S$. Conditional chain rules
retain a vanishing information error at
each fixed $u,t$. The original low-score
list can be unioned over the bounded
parent possibilities described above,
giving a list for each $j$ record.
The mark innovation bounds apply at
$Z_R,\xi_{i,u}$. Radial spreading at
this same parent gives the unnormalized
direction bound of exponent $d$ over
$j$ on the high witness relation.
Finite-prefix likelihood trimming,
with fixed small power slack, verifies
the projection hypotheses at the
actual samples.
This is Lemma~\ref{names:parent-direction-transfer}
with side $Z_R$, auxiliary name $X_F^*$,
$S=(C_*,X_{q_0})$, moving mark $i$ at
parent depth $b=u$, and pin record $j$.
The required separation exponent can be chosen
so that $u-\chi>t$, since $t\ll u$.
Thus $U$ has incremental rate at least
$d$ in the strict-error sense.

This last argument does not assume
homogeneous horizontal increments under
the finer center conditioning. More
explicitly, an expected increment rate
strictly below $2m+d$ along a sequence
$t\downarrow0$ would, by bounded
relative capacities and the two
residual angular rates, give a
positive-probability strict deficit
for the $U$ list. Radial witnesses
provide the required finite-prefix
bounds with probability tending to
one as $t$ decreases. Select one
sufficiently small fixed $t$, take
the tangent information errors to
zero first, and apply the projection
rule to exclude this deficit.
Integration over the open wedge,
followed by Lipschitz endpoint
control, proves \eqref{eq:e-q}.
\end{proof}

\begin{proposition}[Binary time estimate]\label{int:binary}
Assume quadratic admissible data, $m>0$, $c>3$,
and line avoidance. At a binary layer $p\to2p$,
freeze contact coefficients at $F_p$, and let
\[
 b_{\rm miss}(p)=p^{-1}H_\delta
 \bigl(F_{2p}\mid F_{p,2p},\mathcal D,
                      (L_i,L_j)_{2p}\bigr)
\]
in the limiting entropy law. Then
$0\leq b_{\rm miss}\leq1$ and
$e(p)\geq b_{\rm miss}(p)+2m$.
Put $r_0=c\,2^{-\lfloor\log_2c\rfloor}\in[1,2)$.
For the time lower bound, the term
$(c-1)\overline{|G_i|}$ in
\eqref{eq:time-rate} may be replaced by
\begin{equation}
 \frac12\left[(c-r_0)(\bar\alpha-\bar b_{\rm miss})
                  +(r_0-1)(\bar\alpha-1)\right].
 \tag{bin}\label{eq:bin}
\end{equation}
\end{proposition}

\begin{proof}
The center name and two independent horizontal
forms leave only one scalar coordinate; line
avoidance supplies subpower singular-value
bounds. Relative capacities therefore give
$0\leq b_{\rm miss}\leq1$. The joint
increment of $L_i,L_j$ given $F_{p,2p}$
costs at least $2mp$. Indeed $L_i$
annihilates $\ell_i$ exactly and $L_j$
sees it, and conversely. Given
$Q_i,X_p,\mathcal D$, the center bins
to depth $2p$ have bounded lists by
the fiber Taylor estimate. The
horizontal invariance is linear,
so these two comparisons apply
on the whole binary layer. The
conditional chain rule yields
$e(p)\geq b_{\rm miss}(p)+2m$.

On $(\lambda,r_0\lambda)$ use the
microlayer estimate and
$2\E|G_i|\geq\E\alpha-1$.
Tile $[r_0\lambda,c\lambda]$ by
binary intervals $[p,2p]$, where
$p=c2^{-j}\lambda$. On each such
interval,
\[
 2p\leq c\lambda,\qquad
 p\leq(c-1)\lambda.
\]
The output comparison in
Proposition~\ref{int:differential-cost}
then reads $G_i$ to depth $2p$
in the contact basis at $F_p$.
The Taylor remainder is
$O(\delta^{2p})$, and the coefficient
error in \eqref{int:output-coefficients}
times the parent width is no larger.
Freeze the true $(\beta_i,\mathfrak k_i)$
at depth $p$ as output-predicted
auxiliary data before adding marked
conditioning. Subtract the velocity
using the actual child bin time
$\theta'$; this avoids a
$\lambda$-depth error from an
incorrect terminal-time backflow.

The pair $G_i,G_j$ at this accuracy
recovers the three recentered center
readings and $L_i,L_j$, up to
subpower losses. Hence symmetry and
subadditivity lower-bound the
information from a single output
by half the isotropic binary
increment minus $pb_{\rm miss}(p)/2$.
Logarithmic averaging of that
isotropic increment divided by $p$
gives $\bar\alpha$: integrate its
expected derivative on $[p,2p]$
and use dilation invariance.
The binary interval lengths sum
to $(c-r_0)\lambda$, and the
initial interval has length
$(r_0-1)\lambda$. Their sum proves
\eqref{eq:bin}.
\end{proof}

\subsection{Interior action bounds}

We now combine the phase rates into directions of decrease for
the canonical excess. Put $t=1-a$. Along a path with
$dv/dt=k\ge0$, the inequality below gives
$dE(1-t,v(t))/dt\le-2\Gamma$ wherever it applies.
We call $k=0$ freezing: the horizontal resolution is held fixed.
Recall that $c=c'$ is the rate of the active template branch,
so $c=b_1'$ in the linear cases, including a clipped branch.

\begin{proposition}[Admissible local actions]\label{int:actions}
On a valid branch of \eqref{eq:canon}, the actions
listed below satisfy
\begin{equation}
 (\partial_a-k\partial_v)E\geq2\Gamma.
 \tag{surplus}\label{eq:surplus}
\end{equation}
Except where explicitly allowed, assume $m>0$,
radial spreading in the quadratic case, crossing
in the linear case, and line avoidance for
the quadratic $c>1$ estimates.
\begin{enumerate}[label=(\roman*)]
\item Quadratic freezing: $k=0$ is admissible if
$m\geq k_0/(1+c)$ for $0<c<1$, or
$m\geq k_0/2$ for $c\geq1$.
\item Quadratic motion: for $c\geq1$,
$k=(c+1)/4$ is admissible.
\item At quadratic cost $c=1$, without radial
spreading and allowing $m=0$, $k=1/2$ is
admissible when $d\leq1/4$.
\item Linear freezing: $k=0$ is admissible if
$m\geq k_0$. For $c<1$, it suffices more
generally that $m\geq k_0-c\nu$, with
$m>0$ still required.
\item Linear motion with $0<m\leq k_0$:
for $c\geq1$, use $k=(1+c)/2$; for
$c<1$ and $m\geq k_0/(1+c)$, use $k=c$.
\item Linear motion with $2d\leq c<1$:
$k=1$ is admissible for every $m\in[0,1]$,
including zero. Crossing is needed only
when $m>0$.
\end{enumerate}
All assertions are stable under sufficiently
small errors in their hypotheses and inequalities,
with the corresponding additive error in
\eqref{eq:surplus}.
\end{proposition}

\begin{proof}
Write $x,\pi,G,w,\alpha,e$ for the expected
barred rates. Equations~\eqref{eq:time-rate},
\eqref{eq:width-rate}, and \eqref{eq:canon}
reduce the claim to $S\geq0$, where
\begin{equation}\label{int:surplus-algebra}
 \begin{split}
 T&=\min(c,1)(w+x)+(1-c)_+\pi+(c-1)_+G,\\
 S&=T-1-2d-c+k
 \begin{cases}
 5-2\alpha+e,&\text{quadratic},\\
 3-\alpha,&\text{linear}.
 \end{cases}
 \end{split}
\end{equation}
For $c>3$ in the quadratic case we may
replace the last term in $T$ by
\eqref{eq:bin}. The omitted favorable
$\eta$ term is $\eta c$ in the linear
case and $\eta(c-k)$ in the quadratic
case. It is nonnegative for the stated
speeds.

\emph{Quadratic freezing.}
Suppose first $c<1$. Equations~\eqref{eq:modes}
and \eqref{eq:mix-q} give
\begin{equation}\label{int:quadratic-freeze-algebra}
 T\geq l_*+(1+c)n_*+(1-c)a_*+cm
                   +c\max(m+l_*,a_*+\nu).
\end{equation}
If $\nu\leq k_0$, then $n_*=d+\nu$,
$a_*\geq d-c\nu$, and the last maximum
is at least $m+l_*$. Thus
\[
 T\geq(1+c)l_*+2d+2cm+(1+c^2)\nu
       \geq(1+c)l_*+2d+2cm.
\]
The desired bound follows from
$2cm\geq(1+c)(1-l_*)$. If $m\geq k_0$
this is immediate. Otherwise it is
equivalent to $(1+3c)m\geq(1+c)k_0$,
which follows from $m\geq k_0/(1+c)$
and $1+3c\geq(1+c)^2$ for $c\leq1$.
If $\nu\geq k_0$, then $n_*=1$ and
the maximum is at least $a_*+\nu$.
Equation~\eqref{int:quadratic-freeze-algebra}
gives
\[
 T\geq1+c+l_*+cm+a_*+c\nu
       \geq1+c+2d,
\]
because $l_*\geq d$ and $a_*+c\nu\geq d$.

For $c\geq1$, the center and horizontal
bounds give $w\geq2m+2l_*+x$ and
$G\geq2m+x$. Hence
\[
 T\geq2cm+2l_*+(1+c)d.
\]
If $m\leq k_0$, subtracting $1+2d+c$
gives $(1+c)(2m-k_0)\geq0$.
If $m\geq k_0$, substitute $l_*=1$
and use monotonicity in $m$.

\emph{Quadratic motion.}
For $1\leq c\leq3$, use
$2G\geq\alpha-1$, $x\geq d$,
$e\geq2m+d$, and $k=(c+1)/4$.
Cancellation of the $\alpha$ terms gives
\[
 4S\geq(3-c)(1-d)+2(c-1)m\geq0.
\]
For $c>3$, set $b=\bar b_{\rm miss}\in[0,1]$.
The two width bounds imply
$e\geq2m+\max(b,d)$. Substitution of
\eqref{eq:bin} gives
\[
 4S\geq(c+1)(1+\max(b,d)+2m)
         -2(c-r_0)b-2(r_0-1)-4m-4d.
\]
For $b\leq d$ the right side decreases
in $b$, so its minimum occurs at $d$.
For $b\geq d$ it is affine, so only
$b=d,1$ need be checked. The respective
values are
\[
 (c+3-2r_0)(1-d)+2(c-1)m,
 \qquad 4(1-d)+2(c-1)m,
\]
both nonnegative since $r_0<2$.

At $c=1$, using only $e\geq2m$ and
$x\geq0$ gives
\[
 S=\tfrac12-2d-m+x+\tfrac12e
       \geq\tfrac12-2d.
\]
This proves the third assertion,
including $m=0$.

\emph{Linear freezing.}
If $c<1$ and $\nu>0$, Equation~\eqref{eq:mix-l}
gives
$T\geq2l_*+(1+c)n_*+c\nu$. Its excess
over $1+2d+c$ is
\[
 2\min(m,k_0)-(1+c)k_0
       +(1+c)\min(k_0,\nu)+c\nu.
\]
For $\nu\geq k_0$ this is nonnegative.
For $\nu\leq k_0$, if $m\geq k_0$
it is also immediate. Otherwise
$m\geq k_0-c\nu$ lower-bounds it by
$(1-c)k_0+\nu\geq0$.
If $\nu=0$, angular invariance gives
$w\geq\pi$, so
$T\geq\pi+cx\geq2d+(1+c)l_*$.
The threshold gives $l_*=1$.
For $c=1$ use $w\geq x+2d$ and
$x\geq l_*=1$. For $c>1$ use
$w\geq2d+G$, $G\geq m+l_*$,
and $x\geq l_*$. These yield
$T\geq2d+c(m+1)+1\geq1+2d+c$.

\emph{Linear motion.}
For $c\geq1$, put $k=(1+c)/2$ and
use $2G\geq\alpha-1$. Cancellation
gives $S\geq1+x-2d-m$.
Since $m\leq k_0$, $l_*=d+m$;
the bound $x\geq l_*$ leaves $S\geq k_0$.
For $c<1$, put $k=c$. The $\alpha$
terms cancel directly, leaving
\[
 S=-cm+cx+(1-c)\pi+2c-1-2d.
\]
Use $x\geq l_*$ and
$\pi\geq2d+l_*$. The latter follows
from \eqref{eq:double} if $\nu=0$;
if $\nu>0$, \eqref{eq:mix-l} implies it
because $l_*+n_*\geq2d$. Thus
\[
 S\geq(2c-1)k_0+(1-c)m
   \geq\frac{2c^2}{1+c}k_0\geq0.
\]
Finally take $k=1$ and $2d\leq c<1$.
Equation~\eqref{eq:miss} gives
\[
 S\geq1-2d-m+cx.
\]
When $m>0$, crossing and $\ell_j$
give $x\geq m$; when $m=0$ the
same inequality is trivial. Since
$m\leq1$, this lower bound is at
least $c-2d\geq0$.

Every calculation involves finitely
many linear inequalities with bounded
coefficients on a compact parameter
range. Strict-error versions therefore
give the claimed stability; changing
the speed costs the corresponding
Lipschitz error.
\end{proof}
